\documentclass[12 pt, letterpaper]{hmcpset}
\usepackage{hw-preamble}
\setlist[enumerate, 1]{label = (\roman*)}

\name{Jayden Thadani}
\class{MATH 145 --- Topics in Geometry and Topology}
\assignment{Homework 9}
\duedate{2nd April 2025}

\begin{document}

\begin{problem}[45.]
	In this question, we will shwo that the unweighted random walk on $G = \mathbb{Z}$ is recurrent. Here $\mathbb{Z}$ refers to the  $1$-dimensional grid, with a vertex for each integer and an edge between every pair of adjacent integers. We will use $0$ as the centre vertex.
	\begin{enumerate}
		\item The graph $G^{(k)}$ is the restriction of $G$ to the vertex set
			\[
				V^{(k)} = \{ -k, \dots, -1, 0, 1, \dots, k \}
			\]
			and its frontier is $B^{(k)} = \{ -k, k \}$. Solve the Dirichlet problem for $P$ with
			\begin{align*}
				\Delta P(x) & = 0 & 0 < \lvert x \rvert < k \\
				P(x) & = 0 & \lvert x \rvert = k \\
				P(0) & = 1.
			\end{align*}

		\item Calculate the effective resistance $R_{\text{eff}}$ between $0$ and $B^{(k)}$ in two different ways: by using the Dirichlet current, and by using the energy function $\mathcal{E}(P)$. Verify that the answers agree.
		
		\item Show that the unweighted random walk on $\mathbb{Z}$ is recurrent.
	\end{enumerate}
\end{problem}

\pagebreak

\begin{problem}[46.]
	Consider the following network:
	\begin{enumerate}
		\item By setting up a $5$-by-$5$ linear system, or otherwise, solve the Dirichlet problem where the node in the middle is fixed at $1$ and the twenty-four nodes on the outside are fixed at $0$.

		\item An unweighted random walk begins at the centre node. What is the probability that the walker hits the outer edge before returning to the centre?

		\item What is the effective resistance $R_{\text{eff}}$ of the network from the centre node (the source) to the outer nodes (viewed as a single sink) if each edge has resistance $1$?
	\end{enumerate}
\end{problem}

\pagebreak

\begin{problem}[47.]
	Consider potential functions $Q(x)$ defined on the graph of question 46 of the following type. $Q = 1$ on the central vertex, $Q = a$ on the eight vertices surrounding the centre, $Q = b$ on the next sixteen vertices, $Q = 0$ on the outer twenty-four vertices.
	\begin{enumerate}
		\item Calculate the energy $\mathcal{E}(Q)$ as a function of $a, b$.

		\item Determine the values of $a, b$ that minimise the answer in part (i).

		\item Use (ii) to write down a lower bound for $R_{\text{eff}}$, and compare it with 46 (iii).
	\end{enumerate}
\end{problem}

\pagebreak

\begin{problem}[48.]
	In this question you'll show that the unweighted random walk on the 2D grid is recurrent. this graph $G = (V, E)$ has $V = \mathbb{Z}^{2}$ and edges between those pairs of vertices that lie exactly distance $1$ apart. We take $0$ to be the starting vertex of the random walk
	\begin{enumerate}
		\item Sketch the graph $G^{(k)}$ and its frontier $B^{(k)}$ for $k = 0, 1, 2, 3$.

		\item How many vertices in $B^{(k)}$? How many edges between $B^{(k - 1)}$ and $B^{(k)}$?

		\item We define a potential function $Q$ on $G^{(k)}$ as follows
			\[
				Q(x) = \begin{cases}
					1 & x = 0 \\
					1 - a_{1} & x \in B^{(1)} \\
					\vdots \\
					1 - (a_{1} + \dots + a_{k}) & x \in B^{(k)}
				\end{cases}
			\]
			Here $a_{1}, \dots, a_{k}$ are positive real numbers such that $a_{1} + \dots + a_{k} = 1$. This condition ensures that $Q$ is zero on $B^{(k)}$. Write down a formula for $\mathcal{E}(Q)$ as a function of the $a_{i}$.

		\item The optimisation problem
			\[
				\min \{ c_{1} x_{1}^{2} + \dots c_{k} x_{k}^{2} \mid x_{1} + \dots + x_{k} = 1 \}
			\]
			can be shown to have the following solution:
			\[
				x_{i} = \frac{1 / c_{i}}{1 / c_{1} + \dots + 1 / c_{k}}.
			\]

			Using this fact, obtain the smallest possible value of $\mathcal{E}(Q)$ from part (ii). Write down the resulting lower bound for the resistance $R^{(k)}_{\text{eff}}$ between $0$ and the boundary set $B^{(k)}$ in the graph $G^{(k)}$.

		\item Show that the unweighted random walk on $\mathbb{Z}^{2}$ is recurrent
	\end{enumerate}
\end{problem}

\pagebreak

\begin{problem}[49.]
	In this question you'll show that the unweighted random walk on the 3D grid is transient. This graph $G = (V, E)$ has vertex set $\mathbb{Z}^{3}$, and edges between those pairs of vertices that lie exactly distance $1$ apart. We take $0$ to be the starting vertex of the random walk.

	Consider the following current $K$ defined on the edges in the positive octant of the grid. For each vertex $x = (p, q, r)$ with $p, q, r \ge 0$, write $k = p + q + r$. Current will flow from $x$ to the three neighbours in the positive directions as follows:

	\begin{enumerate}
		\item Show that the net flow out of $(0, 0, 0)$ is $1$.
		\item Show that the net flow out of every vertex other than the origin is zero.
		\item How many vertices $(p, q, r)$ in the positive octant satisfy $p + q + r = k$?
		\item Show that each of the three outflow edges at such a vertex has current at most $2 / (k + 2)^{2}$.
		\item Show that
			\[
				\mathcal{E}'(K) < 6 \sum_{k = 0}^{\infty} \frac{1}{(k + 2)^{2}}
			\]
			and deduce that the unweighted random walk on $\mathbb{Z}^{3}$ is transient.
	\end{enumerate}	
\end{problem}

\pagebreak

\begin{problem}[50.]
	Give a quick and easy proof, based on your work in 49 that any locally finite graph that contains $\mathbb{Z}^{3}$ as a subgraph is itself transient
\end{problem}

\pagebreak


\begin{problem}[51.]
	It is well-known that the \emph{consistently biased} random walk on the 1D grid is transient. This walk is defined by the rule
	\begin{align*}
		\mathbb{P}(x_{n + 1} = a + 1 \mid x_{n} = a) & = p \\
		\mathbb{P}(x_{n + 1} = a - 1 \mid x_{n} = a) & = 1 - p.
	\end{align*}
	for some $1 / 2 < p < 1$.
	\begin{enumerate}
		\item Find a system of weights $C : E \to \mathbb{R}$ such that the random walk with these weights is the same as the biased random walk described above.

		\item Construct a unit flow from the origin that has finite energy. Thus, the walk is transient.

		\item Verify that the energy of the flow that you constructed is infinite when calculated in the unweighted 1D grid.
	\end{enumerate}
\end{problem}	

\end{document}
